\cotefile{100}
\scanpage{28}
\mspage{-0.4-}

\uline{Definition} \uline{0.4}: Let $f:T\handwritten{\to}X$ be a $T$-valued point of $X$. Then $\handwritten{\sigma}\circ(1_G\times f)$ is a morphism from $G\underset{S}{\times}T$ to $X$. Define the morphism
\[
\handwritten{\psi}^{G}_{f}:\quad G\underset{S}{\times}T\handwritten{\longrightarrow} X\underset{S}{\times}T
\]
as $(\handwritten{\sigma}\circ(1_G\times f),p_2)$. If no confusion arises, $\handwritten{\psi}^{G}_{f}$ will be shortened to $\handwritten{\psi}_f$. If $f=1_X$, $\handwritten{\psi}_f$ will be denoted $\handwritten{\Psi}$. $\handwritten{\Psi}$ is simply:
\[
(\handwritten{\sigma},p_2):\quad G\underset{S}{\times}X\handwritten{\longrightarrow}X\underset{S}{\times}X\;.
\]

The image of $\handwritten{\psi}_f$ will be denoted $O(f)$ and called the \uline{orbit} of $f$. Now $X\underset{S}{\times}T$, as a scheme over $T$, has a canonical section, namely $(f,1_T)$. Via this, we set up a fibre product defining $S(f)$:
\diagram{\begin{tikzcd}[ampersand replacement=\&]
S(f)\handwritten{\arrow[r]}\handwritten{\arrow[d]} \& T\handwritten{\arrow}[d,"{(f,1_T)}"] \\
G\underset{S}{\times}T\handwritten{\arrow}[r,"{\handwritten{\psi}_f}"] \& X\underset{S}{\times}T.
\end{tikzcd}}
Now $G\underset{S}{\times}T$ is, of course, a group pre-scheme over $T$, and it is not hard to show that $S(f)$ is a subgroup pre-scheme over $T$. If is called the \uline{stabilizer} of $f$.

In case $T=\operatorname{Spec}(k)$, and $f$ is a closed immersion of a point $x\in X$ into $X$, we shall also write $O(x)$ for $O(f)$, and $S(x)$ for $S(f)$. In case $T=\operatorname{Spec}(\handwritten{\Omega})$ and $\handwritten{\Omega}$ is an algebraically closed field, then $f$ is a geometric point of $X$, and it will usually be denoted by a letter $x,y$, etc. Then note that $O(x)$ is a subset of
% archive sheet number: 27
\scanpage{29}

d) \uline{Construction d'un schéma abélien formel}.

\diagram{\begin{tikzcd}[ampersand replacement=\&,column sep=small,row sep=small]
{} \& {} \& \widehat A' \& \widehat A'' \& \widehat A''' \\
M\arrow[d] \& S_0\arrow[l] \& S'_0\arrow[l] \& S''_0\arrow[l,shift left=.55ex]\arrow[l,shift right=.55ex] \& {} \\
T \& \widehat S\arrow[l] \& \widehat S'\arrow[l,dash] \& \widehat S''\arrow[l,shift left=.55ex]\arrow[l,shift right=.55ex] \& \widehat S'''\arrow[l,shift left=1ex]\arrow[l,dash]\arrow[l,shift right=1ex] \\
{} \& S_0 \& S'_0\arrow[l,dash] \& S''_0\arrow[l,dash,shift left=.5ex]\arrow[l,dash,shift right=.5ex] \& S'''_0\arrow[l,dash,shift left=1ex]\arrow[l,dash]\arrow[l,dash,shift right=1ex]
\end{tikzcd}}

Notons d'abord que si $S_0=\operatorname{Spec} k$ ($k$ corps résiduel) alors $\deleted{\unsure{B|A}}S'_0/S_0$ a une section, donc $S'_0\unsure{-|\to}S_0$ est un morphisme à descente effective universelle, et on trouve un $A_0$ sur $\deleted{\unsure{A}}\unsure{S|\Delta}_0$ qui descend $A'_0$.

Commençons par \deleted{prendre} les complétés formels de $S',S'',S''',A',A'',A'''$ \deleted{\unsure{sur}} relativement aux fibres de $s$, on trouve \deleted{\unsure{un|des}} schéma abélien formel $\widehat A'$ sur $\widehat S'$, avec donnée de descente pour $\widehat S'-\widehat S$. On va commencer \unsure{à|:} descendre dans la catégorie des schémas abéliens formels.

On introduit $T=\operatorname{Spf} W(k)$ ($W$ vecteurs de Witt) et $M={}$schéma formel des modules de $A_0$ ($M$ considéré comme schéma formel). Un lemme facile, qui est pratiquement trivial si $S'$ fini sur $S$ [donc $\widehat S',\widehat S''$, \unsure{\&c.|etc.} sont pro-artiniens] dit que la donnée d'une extension \inserted{formelle} $\widehat A'$ de $A'_0=A_{0\,S'_0}$ [resp. $\widehat A''_0$ de $A''_0$.] \&c. équivaut à la donnée d'un $T$-morphisme $\varphi':\widehat S'\to M$ [resp. $\varphi'':\widehat S''\to M$, \deleted{resp. $\widehat S'''-M$}] \&c. La donnée de descente \inserted{sur $\widehat A'$} \unsure{prolongeant} celle que l'on avait sur $A'_0$ équivaut à la condition $\varphi'\widehat p_1=\varphi'\widehat p_2$. Or on constate aussitôt (\unsure{relemme}, trivial encore en \unsure{le|si} $S'$ fini sur $S$, et \unsure{provenant|recourant} du th. de comparaison formel-algébrique en général) que la suite $\widehat S\longleftarrow\widehat S'\mathrel{\substack{\longleftarrow\\[-.6ex]\longleftarrow}}\widehat S''\deleted{\unsure{\leftleftarrows|\Leftarrow}}$ des schémas formels est exacte, donc on trouve $\varphi:\widehat S-M$ t.q. que $\varphi\widehat f=\varphi'$. Cela signifie qu'on a par descente
% archive sheet number: 28
\scanpage{30}
% watermark, sideways, repeated with the lower impression cut by the edge
\margin{MARAIS\quad \unsure{RAIS}}
% Two long oblique pen strokes cancel the typed block.
\deleted{Paris, le 24 Février 1968\\

Cher Grothendieck,\\

J'espère que ta grippe est terminée, et que tu es maintenant complètement rétabli.

Je t'adresse ici l'\deleted{\unsure{ex|e}}Exp. IV, que tu devrais recevoir lun\deleted{\unsure{n|i}}di si la poste marche bien. Alors que je terminais l'appendice (R-R en Géométrie Analytique), je me suis aperçu que tu avais déjà traité la question dans Exp. \unsure{O|0} : il faudra peut-être le vider.

Je joins ma demande de renouvellement pour le CNRS, en te demandant si tu pourrais me l'apostiller. Les années précédentes, Cartan s'en chargeait, mais il estime maintenant que tu es mieux placé que lui pour le faire.

Dès que je pourrai, en principe dans une semaine, je t'enverrai mes prénotes sur Künneth.

Bien à toi,\\}

\handwritten{L Illusie}\\
% archive sheet number: 29
\scanpage{31}

$\widehat A'$ en un schéma formel abélien $\widehat A$.

Le th. de comparaison alg. formel nous ramène à prouver que $A$ est algébrique.

On est ainsi ramené à prouver le

\uline{Th.} Soit $\deleted{\unsure{B}}R$ anneau loc noeth complet, $S=\operatorname{Spec} R$, $\widehat A$ un schéma \inserted{abélien formel} abélien sur $\widehat S=\operatorname{Spf}\unsure{R|B}$. Conditions équivalentes

(i) $\widehat A$ algébri\deleted{\unsure{fiable|sable}}que, i.e. $\exists$ schéma abélien $A$ sur $S$, tel que $\widehat A$ soit isom. à son complété formel \unsure{le long} de $A_s$.

(ii) $\exists$ morphisme propre \uline{surjectif} \deleted{\illegible{1}} $f:S'\to S$, \deleted{\illegible{1}} tel que l'image inverse $\widehat A'$ de $\widehat A$ sur $\widehat S'$ soit algébrique.

(iii) Pour toute \deleted{algèbre finie $B$ sur $A$} quotient \deleted{\unsure{fini}} \inserted{intègre} $B$ de $A$, de dim $1$, si $B'$ est le normalisé de $B$, et $\deleted{\unsure{\widehat B}}$ $\widehat T'=\operatorname{Specf} B'$, \unsure{---} $\widehat A_{\widehat T'}$ est algébrique.

% slanted
\margin{\uline{NB} La démonstration montre qu'on peut se borner \uline{aux} $x$, $x\in S$, $\dim\{x\}=1$, qui ne sont pas contenus dans \deleted{\unsure{tous}} $\bigcup_i T_i$, où les $T_i$ sont \deleted{\unsure{des}} fermés dans $S$ \unsure{et leur|de leur} famille \uline{dénombrable}.}

\deleted{N.B. On peut se borner à prouver le th. si $\dim R=1$, et on \unsure{utilise} \unsure{alors|là} le \unsure{théorème} de \unsure{descente} \unsure{d'effectivité|l'effectivité} de descente fait \unsure{précédemment}.}

\uline{Dém} \quad (i) $\Rightarrow$ (ii) est trivial.

(ii) $\Rightarrow$ (iii) \quad En introduisant $T'=S'\times_S T$, on voit que $\widehat T'\unsure{-|\to}\widehat T$ ``rend $\widehat A_{\widehat T}$ algébrique''. Or on peut trouver, comme $T$ est un trait complet, un \deleted{\unsure{un|autre}} trait $T_1$ fini \deleted{et plat} sur $T$, et un $T$-morphisme $T_1-T'$. Donc $\widehat A_{\widehat T_1}$ est algébrique, donc définit un schéma abélien $A_{T_1}$ sur $T_1$. Ce dernier est muni d'une donnée de descente rel à $T_1-T$ [grâce au th. de comparaison], il suffit \deleted{\unsure{alors}} de voir que la donnée est effective, ce qui est évident, car $T_1-T$ est \inserted{fid} plat fini, et $A_{1}$ est projectif.
% slanted
\margin{\deleted{Pour utiliser \unsure{le} th. de \unsure{N.O.|N.O.S.}, \unsure{la série|le point} \unsure{singulière|singulier} du lieu singulier est \unsure{fermé|fourni}, \unsure{ceci|ce} est \unsure{évidemment} à \unsure{regarder}.}}
% archive sheet number: 30
\scanpage{32}
\mspage{-0.2-}

projective $n$-space \uline{over} \uline{$\handwritten{\mathbb Z}$}, unless, in a particular chapter, all considerations are over a ground field $k$, in which case they will denote affine $n$-space and projective $n$-space over $k$. $\mathrm{PGL}(n)$ will be the projective group acting on $\handwritten{\mathbb P}_n$, $\mathrm{G}\unsure{\mathbb L|\mathrm L}(n)$ the general linear group acting on $\handwritten{\mathbb A}^{n}$, and $\handwritten{\mathbb G}_m=\mathrm{GL}(1)$. The same conventions on the base scheme hold for these group schemes.

(6) $1_x$ is the identity morphism from $X$ to $X$. $p_1$ and $p_2$ are the projections from $X\underset{S}{\times}Y$ to $X$ and $Y$; $p_{12}$, etc. are the projections from $X\underset{S}{\times}Y\underset{S}{\times}Z$ to $X\underset{S}{\times}Y$, etc. If $f:X_1\handwritten{\to}X_2$ and $g:Y_1\handwritten{\to}Y_2$ are $S$-morphisms , then $f\times g:X_1\underset{S}{\times}Y_1\handwritten{\to}X_2\underset{S}{\times}Y_2$ is the product. If $f:X\handwritten{\to}Y_1$ and $g:X\handwritten{\to}Y_2$ are $S$-morphisms, then $(f,g):X\handwritten{\to}Y_1\underset{S}{\times}Y_2$ is the induced $X$-morphism.

(7) If $X$ is a pre-scheme, and $x\in X$, then $\handwritten{\unsure{\varkappa|\kappa}}(x)$ will denote the residue field of $\underline{o}_{x,X}$. If $X$ is a scheme over a field $k$, $X_k$ will denote the set of points $x\in X$ such that $k\handwritten{\overset{\sim}{\to}}\handwritten{\varkappa}(x)$.

\handwritten{\fbox}{\S 1.} \uline{Definition} \uline{0.1}: A group pre-scheme $G/S$ is a morphism $\pi:G\handwritten{\longrightarrow}S$ of pre-schemes, plus $S$-morphisms $\handwritten{\mu}:G\underset{S}{\times}G\handwritten{\to}G$, $\handwritten{\beta}:G\handwritten{\to}G$, $e:S\handwritten{\to}G$ satisfying the usual identities:

(a) \uline{Associativity}:
\diagram{\begin{tikzcd}[ampersand replacement=\&]
G\underset{S}{\times}G\underset{S}{\times}G
\handwritten{\arrow}[r,"{1_G\times\handwritten{\mu}}"]\handwritten{\arrow}[d,"{\handwritten{\mu}\times1_G}"] \& G\underset{S}{\times}G\handwritten{\arrow}[d,"{\handwritten{\mu}}"] \& {} \\
G\underset{S}{\times}G\handwritten{\arrow}[r,"{\handwritten{\mu}}"] \& G \& \text{commutes.}
\end{tikzcd}}
% archive sheet number: 31
\scanpage{33}

\deleted{(iii) $\unsure{\longrightarrow|\Rightarrow}$ (i) (Cas $R$ réduit)}

\[
\text{(iii)}\underset{\text{(ii')}}{\Rightarrow}\left\{\begin{aligned}
&\exists\ \text{morphisme \deleted{\unsure{un}} fini surjectif net}\quad S'\to S,\ \text{tel que}\\
&\widehat A_{\widehat S'}\ \text{soit algébrique.}
\end{aligned}\right.
\]

Soit
\[
\begin{aligned}
&\deleted{\unsure{N_n=P_n|N_n\neq P_n}}\\
&N_n=\underline{NS}_{A_n/S_n}\qquad,\qquad
P_n=\text{\deleted{\unsure{partie}}schéma ouvert induit}\\
&\hspace{15em}\text{défini par les \uline{polarisations}}\\
&\hspace{13em}\text{\deleted{\unsure{ferm}}}
\end{aligned}
\]
Les $P_n$ se recollent et forment un schéma formel $\widehat P$ sur $\operatorname{Spf} R=\widehat S$, \deleted{\unsure{somme}} \inserted{provenant d'une famille \uuline{dénombrable}} de schémas finis nets $P^\alpha_{\deleted{\unsure{n}}}$ sur $S$
\[
\widehat P=\coprod_\alpha\widehat{P^\alpha}
\]

Soit
\[
Z_\alpha=\operatorname{Im}(P^\alpha\to S)
\]
L'hyp. (iii) signifie que pour tt $B$ comme dans l'énoncé, il existe un \inserted{$\alpha$ et un} $S$-morphisme \unsure{$\backslash$|!}
\[
\begin{aligned}
&\deleted{\operatorname{Spec} B'}\\
&T'\longrightarrow P^\alpha
\end{aligned}
\]
donc le pt générique $x$ de $T=\operatorname{Spec} B$ est dans $Z_\alpha$. Donc
\[
Z=\bigcup_\alpha Z_\alpha
\]
contient tous les $x\in S$ tels que $\dim\overline{\{x\}}=1$.

Par un lemme général, il s'ensuit qu'il y a une suite \uline{finie} de $Z_\alpha$ qui recouvrent $S$. Nous prenons pour $S'$ \deleted{\unsure{est}} la somme des $P^\alpha$ correspondants.

(ii') $\Rightarrow$ \deleted{(i) (Cas $R$ réduit)} \inserted{(i') :} $\widehat A_{\widehat S_{\unsure{1|3}}}$ est algébrique, si $S_1=S_{\text{réd}}$.

On peut supposer déjà $S$ réduit. Quitte à faire une extension étale finie $S_1\to S$ (et redescendre, \deleted{\unsure{$\exists$}} ce qui est possible par Raynaud), on peut supposer que
% archive sheet number: 32
\scanpage{34}
\mspage{iii}

a $k$-rational point $x$ is not semi-stable for the action of a reductive group $G$, then there is a one-parameter subgroup $\handwritten{\lambda}$ of $G$, rational over $k$, such that
\[
\handwritten{\mu}(x,\handwritten{\lambda})<0
\]

Roughly, this amounts to asking whether non-semi-stable points $x$ are \uline{so} exceptional that there is some \uline{one} $\handwritten{\lambda}$ (which is therefore rational when $x$ is rational) with respect to which they are most unstable. One approach to this involves the ``Tits fixed point conjecture'' which I would like to mention: if a finite group acts isometrically on a convex set in the flag complex of a semi-simple group, there is a fixed point.

It is certainly not necessary to justify the fact that this book is written entirely in the language of schemes. Of course, the results, for most purposes, could have been stated and proven in a classical language. However, it seems to me \inserted{\handwritten{that it is} only in the language of schemes} that algebraic geometry fulfills
% The circled typed phrase was here, before the second "that". Its arrow points to "me | that".
that essential requirement of all contemporary mathematics: to state its definitions and theorems in their natural abstract and formal setting in which they can be \deleted{considered}\inserted{\handwritten{developed}} independently of geometric intuition. Moreover, it seems to me incorrect to assume that any geometric intuition is lost thereby: for example, the underlying variety in an algebraic scheme is rediscovered, and perhaps better understood through the concept of geometric points. As another example, the theory of schemes has made it possible, in a very intuitive way, to finally dispose of that famous embar\inserted{\handwritten{r}}assment to
% archive sheet number: 33
\scanpage{35}

les extensions résiduelles de $S'\to S$ sont triviales. \deleted{et que les composantes irréductibles de $S$ sont \unsure{géom.} unibranches.} \inserted{Alors} pour \unsure{tte|tt} composante connexe $S'_i$ de $S'$, $S'_i\unsure{-|\to}S$ est une imm. \deleted{\unsure{Si $S$ est}} fermée. Si $S$ est irréductible, une des immersions fermées \deleted{est} un isom (car la somme des images est tout) et on gagne. \deleted{Cela montre déjà que \unsure{pour tte composante|Pour tt composant}} \inserted{On utilise ensuite un}

\uline{Lemme} Soit $\widehat X$ schéma formel propre sur $\widehat S$ ($S=\operatorname{Spec}\deleted{\unsure{A}}R$, $R$ anneau local noeth complet), avec $S$ réduit. Supposons que pour \unsure{tte|tt} composante irréductible $S_i$ de $S$ (avec structure induite réduite) $\widehat X\times_{\widehat S}\widehat S_i$ soit algébrique. Alors $\widehat X$ est algébrique.

Facile par récurrence sur nb de composantes irréductibles, on est ramené au cas où $S=S'\cup S''$, $S'$ et $S''$ sous-schémas fermés, et $\widehat X_{\widehat S'}$ et $\widehat X_{\widehat S''}$ \deleted{\unsure{resp.}} algébriques, provenant de $X'$ et $X''$. Alors $X'_T\simeq X''_T$ ($T=S'\cap S''$) par th. de comparaison, et on recolle $X'$ et $X''$ par cet isom. On voit facilement que ce $X$ recollé marche.

(i') $\Longrightarrow$ (i) \quad Cela résulte du lemme suivant (sans doute cas particulier d'un \deleted{lemme} résultat à démontrer avec complexe cot. relatif ....)

\uline{Lemme} Soit $\widehat X$ schéma formel \inserted{propre et} lisse sur $\widehat S$, où $S=\operatorname{Spec} R$, $R$ anneau local noeth complet. Soit $S'=S_{\text{red}}$, et supposons que $\widehat X_{\widehat S'}$ soit algébrique. Alors $\widehat X$ l'est aussi.
% archive sheet number: 34
\scanpage{36}
\mspage{ii}

some property of semi-simple algebraic groups in all characteristics which takes the place of the full reducibility of representations which is valid only in characteristic $0$. I conjecture, in fact, that if\unsure{'} a semi-simple algebraic group $G$ is represented in a vector space $V$, and if $V_{\unsure{o|0}}$ is an invariant subspace of codimension $1$, then \uline{for} \uline{some} \uline{$\handwritten{\nu}$} the invariant subspace of codimension $1$
\[
V_o\cdot S^{p^{\handwritten{\nu}}-1}(V)\handwritten{\subset}S^{p^{\handwritten{\nu}}}(V)
\]
is complemented by an invariant $1$-dimensional subspace.\textsuperscript{*} Nonetheless, this is unknown, and one consequence is that this book is divided fairly sharply into two halves. Although both parts are closely analogous, they are logically independent. One half consists in Chapters 1, 2, 4, and 5 which deal essentially, only with characteristic zero, and yield a construction of the moduli scheme for curves over $\handwritten{\mathbb Q}$. The other half consists in Chapters 3, 6, and 7, which deal with the ``arithmetic \unsure{c|\v{c}}ase'', i.e. over $\operatorname{Spec}(\handwritten{\mathbb Z})$, and yield a construction of the moduli schemes for curves over $\handwritten{\mathbb Z}$. From a standpoint of content, however, Chapters 1, 2, 3, and 4 deal with orbit space problems, while Chapters 5, 6, and 7 deal with moduli.

Another unsolved problem deals with rationality over a perfect but non-algebraically closed field $k$, and is related to the Godement conjecture. In the notation of Chapter 2, the problem is whether, when

\noindent\rule{5em}{0.4pt}\par
\noindent\textsuperscript{*}Here $S^k V$ stands for the $k^{\mathrm{th}}$ symmetric power of $V$.
% archive sheet number: 35
\scanpage{37}

On est ramené à prouver que si \deleted{\unsure{$S'$ est}} on a un sous-schéma $S_0=V(J)$, $J^{\unsure{2|n}}=0$, et si $\widehat X_{\widehat S_{\unsure{0'|0}}}$ est algébrique, alors $\widehat X$ l'est aussi. On utilise la théorie du prolongement infinitésimal d'une structure lisse $X_0/S_0$ resp $\widehat X/\widehat S_0$, et :

$1^\circ$) L'\deleted{\unsure{déf}}\unsure{obstruction} à prolonger \deleted{\unsure{S}}$X_0$ (resp $\widehat X_0$) en un $X$ (resp. $\widehat X$) lisse sur $S$ (resp $\widehat S$) est dans ``le $=$'' groupe
\[
H^2(X_0,\unsure{\mathscr G|\mathscr T|\mathfrak g|\mathcal T}_{X_0/S_0}\otimes_{\mathcal O_{S_0}}J)\simeq H^2(\widehat X_0,\mathscr G_{\widehat X_0/\widehat S_0}\widehat\otimes_{\mathcal O_{\widehat S_0}}\widehat J
\]
et moyennant cette identification, l'obstruction est la $=$.

$2^\circ$) L'indétermination pour remonter est dans ``le $=$'' groupe
\[
H^1(X_0,\mathscr G_{X_0/S_0}\otimes_{\mathcal O_{S_0}}J)\simeq H^1(\widehat X_0,\widehat{\mathscr G}_{\widehat X_0/\widehat S_0}\widehat\otimes_{\mathcal O_{\widehat S_0}}\widehat J).
\]

Il reste à prouver la

\uline{Proposition} Soit $\deleted{\unsure{R}}S$ un \deleted{anneau local noeth complet.} \inserted{schéma local \inserted{noeth} complet}, \deleted{de dim $\unsure{>| \geq}1$}, $(Z_i)_{i\in\deleted{\unsure{N|\mathbb N}}I}$ une \deleted{suite} \inserted{famille dénombrable} de parties fermées \deleted{\unsure{non}} de $S$, telles que $Z=\bigcup_i Z_i$ contienne tt $x\in S$ avec $\dim\overline{\{x\}}\leq1$. Alors il existe une partie finie $J$ de $I$ telle que $\bigcup_{i\in J} Z_i=S$.

\inserted{Soit $S=\operatorname{Spec} R$.} Si $\dim\deleted{\unsure{B|S}}\unsure{S|R}\leq1$, c'est trivial (alors \inserted{l'ens sous-jacent à} $S$ est fini, et l'hyp. signifie que $Z=S$). Si $n\mathrel{\unsure{:=|:|=}}\dim S\geq\unsure{2|?}$, on sait qu'il existe un anneau à val. discrète complet $V$, \deleted{tel} tel que $R$ soit fini sur $V[[T_1,\ldots,T_{\unsure{r|n}}]]$, \inserted{$\unsure{r=n-1}$}. Notons qu'on est ramené \inserted{aussitôt} à prouver
% archive sheet number: 36
\scanpage{38}

\begin{center}
\uline{Preface}
\end{center}

The purpose of this book is to study two related problems: when does an orbit space of an algebraic scheme acted on by an algebraic group exist? And to construct moduli schemes for various types of algebraic objects. The second problem appears to be, in essence, a special and highly non-trivial case of the first problem. From an Italian point of view, the crux of both problems is in passing from a birational to a biregular point of view. To construct both orbit spaces and moduli ``generically'' are simple exercises. The problem is whether, within the set of all models of the resulting birational class, there is one model whose geometric points classify the set of orbits in some action, or the set of algebraic objects in some moduli problem. In both cases, it is quite possible that some orbits, or some objects are so exceptional, or, as we shall say, are \uline{not} \uline{stable}, so that they must be left out of the model. The difficulty is to pin down the meaning of stability in a given case. One of the most intriguing unsolved problems, in this regard, is that of the moduli of non-singular polarized surfaces. Which such surfaces are not stable, in the sense that there is no moduli scheme for them and their deformations? This property is very delicate.

One of my principles has been not to worry too much about the difference between characteristic $0$, and finite characteristics. A large part of this book is, therefore, devoted to a theory developed only in characteristic $0$. I am convinced, however, that it is almost entirely valid in all characteristics. What is necessary is to find
% archive sheet number: 37
\scanpage{39}

le

\uline{Corollaire} Supposons que les $Z_i$ soient \uline{\unsure{rares}} dans $S$. Alors il existe un $x\in S$, avec $\dim\overline{\{x\}}=1$, tel que $x\notin\bigcup Z_i$.

Remplaçant $S$ par $T=\operatorname{Spec} V[[T_1,\ldots,T_{\unsure{r|n}}]]$, \deleted{\unsure{on a}} on voit que les $T_i:=\operatorname{Im}(Z_{\unsure{i}}\unsure{-|\to}T)$ sont rares, et il suffit de trouver un $x\in T$, avec $\dim\overline{\{x\}}=1$, tel que $x\notin\bigcup T_i$. (En effet, il existe alors un $y\in S_{\unsure{2|x}}$ au dessus de $x$, et on a nécessairement $\dim\overline{\{y\}}=1$).

On va montrer qu'il existe $\unsure{\overline{\mathrm{m}}|\overline{\mathrm{me}}}$ un $x\in T$ provenant d'un $V$-hom
\[
V[[T_1,\ldots,T_r]]\unsure{-|\longrightarrow}V\;.\qquad \deleted{\text{\unsure{Comme}}}\unsure{\text{,}}
\]

En effet, l'ens. des ces hom s'identifie à $\mathfrak m^{\unsure{r|n}}$ ($\mathfrak m=\operatorname{rad}(\unsure{V|W})$). Il a une structure d'espace métrique complet. Montrons que pour tte partie fermée rare $Z$ de $T$, l'ens des $\varphi$ tels que $V(\operatorname{Ker}\varphi)\subset Z$ est rare. En effet, on peut supposer que \deleted{$Z$ est \unsure{défi|défini}} $Z=V(f)$, \inserted{$f\in V[[T_1\unsure{,|-}\ldots T_r]]$, $f\neq0$,} et $\lvert V(\operatorname{Ker}\varphi)\rvert\subset\lvert V(f)\rvert$ équivaut (comme $V(\operatorname{Ker}\varphi)$ est intègre \unsure{bien|donc} évidemment) à $V(\operatorname{Ker}\varphi)\subset V(f)$ \unsure{i.e.|\ldots} $f\in\operatorname{Ker}\varphi$ i.e. $\varphi(f)=0$; i.e. $f(\varphi_1,\ldots,\varphi_r)=0$. Il est immédiat que c'est une partie fermée rare.
% archive sheet number: 38
\scanpage{40}
% Three oblique pen strokes cross the typed letter.
\deleted{7.4.1967\\

Cher Dieudonné,\\

D'accord qu\deleted{\unsure{i|l}}'il est inutile de rajouter un (vi bis) en Errata à 12.1.1. ; il y aurait lieu de le faire cependant en réédition, car le lecteur trouvera plus faci\deleted{\unsure{el|le}}lement dans 12.1.1. (où il regardera tout naturellement) que dans 12.1.3. I\deleted{\unsure{o|a}}l n'y a pas de raison d'avoir mis l'énoncé sur la prof. rectifiée à part des autres. Certainement tu te canules dans 15.5.6.1. car il est clair à priori que si l'énoncé est vra\unsure{i} lorsque tout fermé irréductible a un \deleted{seul} point irréductible unique, il est vrai aussi sans \deleted{\unsure{ce|de}}la restriction d'unicité: il suffit de regarder l'espace \inserted{$X'$} formé des parties fermées irréductibles et l'application $\deleted{\unsure{E}}X\quad X'$ qui induit un isomorphisme sur les \deleted{\unsure{les}} ensembles ordonnés des ouverts de $X$ et $X'$ ... Il me semble certain par suite qu'en changeant une ligne ou deux dans ta démonstration, elle doit marcher dans le cas général. Dans 15.6.3. on peut éviter l'hypothèse noethérienne, car la réduction au cas d'un anneau de valuation discrète est bien inutile, on se ramène à vérifier le relèvement des générisations, qui résu\deleted{\unsure{o|a}}lte de 15.6.2. D\deleted{,}\inserted{'}accord par contre que la démonstration donnée de 15.6.4. utilise effectivement l'hypothèse noethérienne; on l'él\deleted{\unsure{é|m}}iminera dans la réédition, en faisant une réduction au cas noethérien.

As-tu fait le nécessaire pour qu'on m'envoie la suite des 2èmes\unsure{:|.} épreuves ?

Bien cordialement\\}
% archive sheet number: 39
